% 本文件由 LaTeX 排版 Agent 根据有效 translation.json 自动生成。
% author=Ambrose; work_id=3405; title=论奥迹

\chapter{论奥迹}

% structure: p0001 source=h1 target=omitted reason=page-title-already-rendered-by-parent

% structure: p0002 source=h2 target=section reason=relative-heading-rank; base=h2

\section{引言}

著作者在本论著开头说明，他的目的乃是为那些即将领洗的人，阐述那件圣事、以及坚振与圣体圣事的仪式和意义。因为在早期教会中，这些事都是极度保密的，惟恐被外教人亵渎，而且按照习惯，正如著名的耶路撒冷的\Person{圣济利禄}的\WorkTitle{教理讲授}那样，都是在四旬期后半段向慕道者讲解这些事。\par

因此，这类论著具有特殊的价值，因为我们在其中可以清楚地看到，在那些流传至今的讲道词被编纂的时代，教会的完整训导。\par

\Person{圣安博}逐一讲解并阐释了大部分内容：首先是庄严圣洗时常用的仪式，指出这些外在事物所蕴含的深奥真理与奥迹；然后他论述坚振圣事，提及圣神的七种恩赐；最后，他讲述圣体圣事，特别阐明了真实临在的教义。\par

自十六世纪以来，有些作家试图证明这本论著并非\Person{圣安博}所作，但此事毫无真正的疑问，正如本笃会编辑者们所确证、如今已普遍承认的那样。该论著是为四旬期使用而编写的，但无法确定具体年份，可能是在公元387年左右，根据对\WorkTitle{论圣祖}一书的引用推断。\par

% structure: p0007 source=h2 target=section reason=relative-heading-rank; base=h2

\section{第一章}

\Person{圣安博}说，在他已经讲解了圣善生活之后，现在要讲解奥迹。然后，在说明了自己之前为何没有这样做的理由之后，他解释了开耳奥迹，并指出这早先是由基督亲自施行的。\par

1. 我们每日都在谈论有关伦理的主题，当诵读圣祖的事迹或\BibleBook{}{箴言}的训诲时，为的是使你们通过这些教导和训示，逐渐习惯于踏上古人的道路、行走在他们的途径上，并遵从天主的诫命；为的是使你们因圣洗而更新后，能保持那与受过洗者相称的生活方式。\par

2. 现今的节期提醒我们该讲述奥迹了，并阐明诸圣事的含义；如果我们认为在领洗之前，就该将这些奥迹教导给尚未入门的人，那我们就可被视为是泄漏了奥迹，而非描绘了奥迹。此外，还有一个理由，就是奥迹本身的光芒，对那些正在期待而不知究竟的人，会比事先有过任何讲解更有效地照耀他们。\par

3. 所以，敞开你们的耳朵吧，吸入由圣事的恩宠所吹拂于你们的永生芳香；这正是我们在举行开耳奥迹时向你们所表示的，当时我们说：\BibleQuote{厄法达，就是说：开了吧！} \BibleRef{谷 7:34}，好使那前来寻求平安的人，能知道他所被问的是什么，并能记住他所回答的。\par

4. 我们在\BibleBook{}{福音}中读到，基督在治愈那又聋又哑的人时，就用了这奥迹。但祂触摸了他的口，因为那被治愈者是哑巴，而且是男性；就一方面而言，是要他开口发出所赋予的声音；就另一方面而言，是因为那触摸对于男性是合宜的，而对于女性则不合宜。\par

% structure: p0013 source=h2 target=section reason=relative-heading-rank; base=h2

\section{第二章}

论及那些即将入门者在进入教会时应许了什么，这些许诺的见证者，以及他们为何随后转向东方。\par

5. 此后，至圣所就为你们敞开了，你们进入了重生的圣所；回想你们被问了什么，并记住你们回答了什么。你们弃绝了魔鬼及其作为，弃绝了世俗及其奢华和享乐。你们的那句话，不是保存在死者的墓穴里，而是保存在生命册上。\par

6. 你们在那里看见了执事，看见了司铎，还看见了主祭[即主教]。不要考虑外在的形体，而要思量奥迹的恩宠。你们是在天使面前说话，正如经上所载：\BibleQuote{因为司祭的唇舌应保卫智识，人也应从他的口中得到教训，因为他是万军上主的使者。} \BibleRef{拉 2:7} 在此没有欺骗或否认的余地。他是天使，宣讲基督的国度和永生。你们不应该按他的外表，而应按他的职分来尊重他。想一想他所传授的，思考一下他给予你们的生活准则，认清他的地位。\par

7. 因此，你们进来，为能认清你们的仇敌，你们就是当着他的面弃绝了他；然后你们转向东方；因为凡弃绝魔鬼的，就是转向基督，且面对面地注视祂。\par

% structure: p0018 source=h2 target=section reason=relative-heading-rank; base=h2

\section{第三章}

\Person{圣安博}指出，我们必须考虑那水中和圣职者身上的天主性临在与作为，然后他引用了许多旧约中圣洗的预象。\par

8. 你们看见了什么？水，的确，但不仅仅是水；你们看见执事们在那里服务，主教在提问并祝圣。首先，宗徒曾教导你们，不该注目于\BibleQuote{那看得见的，而该注目于那看不见的；那看得见的，原是暂时的；那看不见的，才是永远的。} \BibleRef{格前 5:18} 因为你们在别处读到：\BibleQuote{其实，自从天主创世以来，祂那看不见的美善，即祂永远的大能和祂为神的本性，都可凭祂所造的万物，辨认洞察出来。} \BibleRef{罗 1:20} 因此，主自己也说：\BibleQuote{纵使你们不肯信我，至少要信这些工作。} \BibleRef{若 10:38} 那么，就请相信，天主性的临在就在那里。你们相信这作为，却不相信这临在吗？除非临在先存，否则作为又从何而来呢？\par

9. 然而，请考虑这奥迹是多么古老，甚至在世界的创造之初就已被预像。在最起初，当天主创造天地时，经上说：\BibleQuote{圣神，运行在水面上。} \BibleRef{创 1:2} 那运行在水面上的，难道不是在水面上工作吗？但我为什么说\Quote{工作}呢？就祂的临在而言，祂在运行。那运行的难道不是在工作吗？你要认识到，在世界的创造中，祂曾工作，正如先知所说：\BibleQuote{因天主的一句话，诸天造成；因祂口中的圣神，万象生成。} 每一项陈述都基于先知的见证：既是祂在运行，也是祂在工作。\Person{梅瑟}说祂运行了，\Person{达味}则作证祂工作了。\par

10. 再取另一见证。所有血肉因邪恶都已败坏。天主说：\BibleQuote{我的圣神将不永存于人中，因为他们都是血肉。} \BibleRef{创 6:3} 天主借此表明，圣神的恩宠因肉性的不洁和重罪的污秽而离去。因此，天主愿意恢复所亏损的，就使洪水泛滥，并命义人\Person{诺厄}进入方舟。及至洪水退去，诺厄先放出一只不归的乌鸦，又放出一只鸽子，据说它衔着一枝橄榄叶回来了。你看到水，你看到木头（方舟的木），你看到鸽子，难道还对这奥迹犹豫吗？\par

11. 因此，水是那肉身入浸其中的，为洗净一切肉性的罪。一切邪恶都埋葬于此。木头是主耶稣为我们受苦时，被钉于其上的那木。鸽子是圣神降下的形像，正如你在新约中所读到的，祂赋予你灵魂的平安和心灵的宁静。乌鸦是罪的预像，它飞去不再归来，如果你内外都保持着义德的话。\par

12. 还有第三个见证，就如宗徒教导我们的：\BibleQuote{我们的祖先都曾在云柱下，都从海中走过，都曾在云中和海中受了洗归于\Person{梅瑟}。} \BibleRef{格前 10:1} 此外，梅瑟自己在凯旋歌中说：\BibleQuote{你吹了你的圣神，海水就将他们淹没。} \BibleRef{出 15:10} 你观察到，在那希伯来人过海的事件中，圣洗圣事早已被预表：埃及人灭亡，而希伯来人生还。我们每天在这圣事中所受的教导，岂不是罪过被吞灭、错误被消除，而德行与无辜得以保全无伤吗？\par

13. 你听到我们的祖先曾在云柱下，那是一柱仁慈的云，熄灭了肉欲的燥热。那仁慈的云荫庇着那些圣神所造访的人。最终，它临于\Person{童贞玛利亚}，至高者的德能荫庇了她，\BibleRef{路 1:35} 当时她为人类孕育了救赎。而这奇迹通过\Person{梅瑟}预像性地完成了。那么，如果圣神曾在预像中，祂在现实中岂不更临在吗？因为圣经对我们说：\BibleQuote{法律是由\Person{梅瑟}传授的，但恩宠和真理却由\Person{耶稣基督}而来。} \BibleRef{若 1:17}\par

14. \Place{玛辣}是一处极苦的水泉：\Person{梅瑟}投入木头，水就变甜了。因为水若没有主十字架的宣讲，于未来的救恩毫无裨益，但经过救赎十字架奥迹的祝圣后，它就适合用作灵性之洗和救恩之杯。因此，正如梅瑟——即那位先知，将木头投入那泉中，同样，司铎在这圣洗池上宣讲主十字架的讯息，水便为施予恩宠而变得甘甜。\par

15. 那么，你切不可完全信赖你肉身的眼目；那看不见的才更真实地被看见，因为所见的是暂时的，而另一者却是永恒的，非肉眼所及，而是由心灵和神魂辨明的。\par

16. 最后，让最近从\WorkTitle{列王纪}中所读的训诲教导你吧。\Person{纳阿曼}是叙利亚人，患有癞病，无人能使他洁净。后来，有一个被掳的少女说，在\Place{以色列}有一位先知，能治好他的癞病污秽。据说，他带了金银去见以色列王。王听见他来的原因后，便撕裂自己的衣服，说，这乃是在寻隙攻击他，因为向他要求的事，并非君王权能所及。然而，\Person{厄里叟}传话给君王，叫他将那叙利亚人送到自己那里，好让他知道以色列中有天主。及至他来，便命令他在\Place{约旦}河里浸洗七次。\par

17. 于是，他心中寻思，自己家乡有更好的水，他曾多次在其中沐浴，却从未从癞病得洁净；因此记起这事，便不肯听从先知的命令，但后来因仆人们的劝谏与说服，他顺从了，将自己浸入水中。他立刻得了洁净，便明白了，使人洁净的并非水，而是恩宠。\par

18. 现在你该明了，那被掳中的少女是谁。她就是由外邦人中聚集的会众，也就是\Emph{天主的教会}，昔日被罪恶的俘虏所奴役，尚未享有恩宠的自由，借着她的劝导，那愚蠢的外邦人民曾听闻了先前所怀疑的预言圣言。后来，当他们相信应当听从时，便被洗净了所有罪的污秽。那个癞病人确实在未得治愈前怀疑过；你已经得了治愈，因此不该再怀疑。\par

% structure: p0031 source=h2 target=section reason=relative-heading-rank; base=h2

\section{第四章}

水若没有圣神则不能洁净，这由\Person{若翰}的见证和圣事施行形式本身就显示出来。并且，这也由福音中的水池及其被治愈之人所预表。在同一段经文中，还显明圣神确实在\Person{基督}受洗时降临于祂身上，并解释了这一奥迹的意义。\par

19. 先前告诉你们不要只相信所见之事的原因，是为了让你们或许不致说：这就是那伟大的\OriginalTerm{奥迹}{mysterium}，\BibleQuote{眼所未见，耳所未闻，人心所未想到的}。\BibleRef{格前 2:9} 我看见了水，那是我每天惯常见到的水。难道这我常常在其中沐浴却从未得到洁净的水，现在就要洁净我吗？由此你可以认出，没有圣神，水不能洁净。\par

20. 因此，你们要读一读：圣洗中的三个证据——水、血和圣神——\BibleRef{若一 5:7} 是一致的，因为若除去其中任何一个，圣洗圣事就不存在了。因为没有基督十字架的水算什么呢？那是普通的元素，没有任何圣事的效力。同样，没有水，就没有重生的圣事：\BibleQuote{人除非由水和圣神重生，不能进入天主的国。}\BibleRef{若 3:5} 然而，就连慕道者也相信主耶稣的十字架，他也因此被划上十字圣号；但除非他因父及子及圣神之名受洗，就不能获得罪过的赦免，也不能领受灵性恩宠的恩赐。\par

21. 因此，那个叙利亚人在法律下浸入水中七次，而你们却因天主圣三的名受洗，你们宣认了圣父。请回想你们所做的：你们宣认了圣子，宣认了圣神。要清楚这信仰中的次序：你们向世界死了，并向天主复活了。并且，你们在那水中仿佛被埋葬于世界，对罪死了，却复活得永生。所以，你们要相信，这些水并非毫无能力。\par

22. 因此，经上说：\BibleQuote{主的天使按时下到水池中，搅动池水；水被搅动后，第一个下去的人，无论患什么病，必得痊愈。}\BibleRef{若 5:4} 这水池在\Place{耶路撒冷}，每年有一人得痊愈，但在天使降下之前，无人得痊愈。由于那些不信的人，池水被搅动，作为天使下降的标记。他们有一个标记，你们却有信德；为他们有一位天使降下，为你们却有圣神；为他们受造之物被搅动，为你们却是基督自己，受造之物的主，在行事。\par

23. 那时，只有一人得痊愈，现在所有人都得痊愈；或更确切地说，只有基督徒子民得痊愈，因为在某些人那里，连水也是欺诈的。\BibleRef{耶 15:18} 不信者的圣洗不能治愈，反而玷污。犹太人洗濯杯盘，好像无灵之物也能承担罪过或蒙受恩宠。但你要洗濯你这活的杯爵，好让你的善工在其中闪耀，你恩宠的光荣得以辉耀。因为那水池是一个预像，为使你相信天主的能力降临在这圣洗池上。\par

24. 最后，那瘫子一直等候一个人。那人是谁呢？除非是童贞女所生的主耶稣，在祂来临后，不再有影子逐一医治人，而是真理要医治所有人。所以，祂就是那位人等候降临的，圣父曾对\Person{洗者若翰}论及祂说：\BibleQuote{你看见圣神降下，停在谁身上，谁就是那以圣神施洗的。}\BibleRef{若 1:33} 若翰也为祂作证说：\BibleQuote{我看见圣神仿佛鸽子从天降下，停在祂身上。}\BibleRef{若 1:32} 圣神为什么仿佛鸽子降下？正是为了让你看见，让你认出，那\Person{诺厄}从方舟放出的鸽子正是这只鸽子的预像，为使你认出这圣事的预像。\par

25. 或许你会反驳说：既然放出的是一只真实的鸽子，而圣神仿佛鸽子降下，我们怎能说那里是预像，这里才是实体，而在希腊原文中却记载圣神以鸽子的形像降下？但有什么比永远长存的天主性更真实呢？受造物不能是实体，只能是形像，且易于毁灭和改变。同样，因为受洗者的纯朴不应仅是外在的形像，而应是内在的实体，且主说：\BibleQuote{机警如同蛇，纯朴如同鸽子。}\BibleRef{玛 10:16} 所以，祂恰当地仿佛鸽子降下，为劝勉我们应当具有鸽子的纯朴。此外，我们又读到，形像被当作实体，关于基督：\BibleQuote{祂取了人的形像；}\BibleRef{斐 2:8} 关于天主圣父：\BibleQuote{你们从未见过祂的形像。}\BibleRef{若 5:37}\par

% structure: p0040 source=h2 target=section reason=relative-heading-rank; base=h2

\section{第五章}

基督自己临在于圣洗圣事中，因此我们不必考虑施行圣洗者的个人身份。下文简要解释即将受洗者通常所宣认的天主圣三的信德宣认。\par

26. 那么，这里还有任何怀疑的余地吗？既然在福音叙述中，圣父从天上清晰宣告说：\BibleQuote{这是我的爱子，我所喜悦的。}\BibleRef{玛 3:17} 圣子也说话，圣神曾以鸽子的形像显现并停留在祂身上；圣神也说话，祂曾以鸽子的形像降下；连\Person{达味}也说：\BibleQuote{主的声音响彻水面，荣耀的天主雷鸣，主临于大水之上。}圣经纪载：因\Person{耶鲁巴耳}的祈祷，有火从天降下\BibleRef{民 6:21}，同样，当\Person{厄里亚}祈祷时，有火降下并圣化了祭献。\par

27. 不要看个人的功劳，而要看司铎的职务。或者，如果你要看功劳，就要把司铎看作\Person{厄里亚}。也要看\Person{伯多禄}和\Person{保禄}的功劳，他们从主耶稣领受了这奥迹，并传给了我们。对于旧约时代的人，天主降下可见的火，好使他们相信；对于我们这些相信的人，主则无形地工作；发生在他们身上的事，是作为预像；为我们，则是作为警戒。所以，你要相信，主耶稣临在于司铎的呼求中，祂曾说：\BibleQuote{哪里有两三个人，因我的名字聚在一起，我就在他们中间。}\BibleRef{玛 18:20} 那么，在教会所在，在祂的奥迹所在的地方，祂岂不更乐意赐予自己的临在！\par

28. 你然后下降（进入水中），回想你如何回答那些问题：你信圣父，你信圣子，你信圣神。那陈述不是：\Quote{我信一个较大的，一个较小的，和一个最低的位格}，而是你受你口中所发的同样保证之约束，要以同样的方式信圣子，如同你信圣父；并以同样的方式信圣神，如同你信圣子，只有这一点除外，即你宣认你必须唯独信主\Person{耶稣}的十字架。\par

% structure: p0045 source=h2 target=section reason=relative-heading-rank; base=h2

\section{第六章}

为何他们从圣洗池中上来时，头上被傅油；为何，在圣洗之后，他们的脚被洗，以及每一样分别赦免什么罪。\par

29. 此后，你走到司铎那里，想想接下来发生了什么。岂不是\Person{达味}所说的：\BibleQuote{香膏流在头上，流到胡须，流到\Person{亚郎}的胡须}？这就是\Person{撒罗满}也提到的那种香膏：\BibleQuote{你的名香膏倾流，因此众童女都爱你，吸引你。}~\BibleRef{歌 1:2} 主\Person{耶稣}，今日多少重生的灵魂爱了祢，并说：\BibleQuote{求你吸引我，我们将在你身后奔跑，为了你香膏的芬芳，}~\BibleRef{歌 1:3} 好使他们能吸入你复活的芬芳。\par

30. 现在思想这事为何如此行，因为 \BibleQuote{智慧人的眼睛在他头上；}~\BibleRef{训 2:14} 因此香膏流到胡须，就是说，流到青春的美丽；因此，流到\Person{亚郎}的胡须，好使我们也能成为被选的种族，司祭的、宝贵的，因为我们全都被属神的恩宠所傅油，为分享天主的国和司祭职。\par

31. 你从圣洗池上来；记得福音的记载。因为我们的主\Person{耶稣基督}在福音中洗了祂门徒的脚。当祂来到\Person{西满伯多禄}面前时，\Person{伯多禄}说：\BibleQuote{你永不可洗我的脚。}~\BibleRef{若 13:8} 他不明白这奥迹，因此拒绝了这项服事，因为他认为如果容忍主服侍他，仆人的谦逊就会受损。主回答他说：\BibleQuote{我若不洗你的脚，你就与我无分。} \Person{伯多禄}听了，回答说：\BibleQuote{主，不但我的脚，而且我的手和头。} 主回答说：\BibleQuote{沐浴过的人，只需要洗脚，全身就洁净了。}~\BibleRef{若 13:9}\par

32. \Person{伯多禄}是洁净的，但他必须洗脚，因为他从原祖那里继承而有了罪，那时蛇使他跌倒，并诱使他犯罪。因此他的脚被洗，为使原罪得以除去，因为我们自己的罪是通过圣洗赦免的。\par

33. 同时注意，这奥迹正包含在谦逊的职务本身中，因为\Person{基督}说：\BibleQuote{我，你们的主和师傅，尚且洗你们的脚；你们也当彼此洗脚。} 因为，既然救恩的创始者亲自借着服从救赎了我们，我们作祂仆人的，岂不更应献上我们谦逊和服从的服事吗？\par

% structure: p0052 source=h2 target=section reason=relative-heading-rank; base=h2

\section{第七章}

慕道者的白衣标志着罪过的洗涤，因此教会说自己黝黑而秀美。天使惊奇于她的光辉，如同惊奇于主肉身的光辉。而且，\Person{基督}以许多形象向他的新娘称赞自己的美丽。描述了彼此的爱慕。\par

34. 此后，给你们穿上白衣，作为你们脱去罪过的覆盖，穿上纯洁的贞洁细麻衣的标记；对此先知说：\BibleQuote{你要以牛膝草洒我，我就洁净；你要洗涤我，我就比雪更白。} 因为受洗的人，按法律和福音看来，被视为洁净了：按法律，因为\Person{梅瑟}用一束牛膝草蘸羔羊的血洒；~\BibleRef{出 12:22} 按福音，因为当\Person{基督}在福音中显示祂复活的光荣时，祂的衣裳洁白如雪。因此，罪过得赦免的人，变得比雪更白。所以天主借\Person{依撒意亚}说：\BibleQuote{你们的罪虽似朱红，我要使它们变成雪白。}~\BibleRef{依 1:18}\par

35. 教会借着重生的圣洗穿上这些衣服后，在\BibleBook{}{雅歌}中说：\BibleQuote{我虽黝黑，却秀美，\Place{耶路撒冷}的女子啊。}~\BibleRef{歌 1:4} 黝黑是由于人性状况的脆弱，秀美是由于信德的圣事。\Place{耶路撒冷}的女子看见这些衣服，惊奇地说：\BibleQuote{那来的是谁，上升如白衣？}~\BibleRef{歌 8:5} 她曾是黝黑的，如今怎忽然变白？\par

36. 当\Person{基督}复活时，天使也惊奇；天上的掌权者看见肉身升天时，也惊奇。那时他们说：\BibleQuote{谁是这位光荣的君王？} 当有人说：\BibleQuote{王侯，请提高你们的门，永恒的门，请高抬，光荣的主要进来。} 在\BibleBook{依撒意亚}{书}，我们也发现天上的掌权者惊奇说：\BibleQuote{那从\Place{厄东}而来，身穿染红了的服装，那由\Place{波责辣}而来，穿着华美，前进步武的是谁？}~\BibleRef{依 63:1}\par

37. 但是\Person{基督}，看见祂的教会——为这教会，正如你在\BibleBook{匝加利亚}{}先知书中看到的，祂曾穿上污秽的衣服——现在却穿上白衣；就是说，看到一个纯洁的、在重生的圣洗中洗净的灵魂，祂说：\BibleQuote{看，我的爱卿，你多么美丽！看，你多么美丽！你的眼有如鸽子的眼。}~\BibleRef{歌 4:1} 圣神曾以鸽子的形状从天降下。眼美如鸽，因为圣神以鸽子的形状从天降下。\par

38. 接着又说：\BibleQuote{你的牙齿像一群剪了毛的母羊，刚由水池中上来，每只都怀有双胎，没有一只丧失子嗣的；你的嘴唇像一缕朱红线。}\BibleRef{歌 4:2}这不失为重大的赞美。首先藉那剪毛之羊的可爱比喻；因为我们知道，山羊既能在高处吃草而无危险，在崎岖之地也能安全觅食，剪毛之后又除去多余之物。教会正可比作这样一群羊，她内里拥有许多灵魂的美德，它们通过在\OriginalTerm{圣洗池}{laver}中涤除了多余的罪过，并向基督献上\OriginalTerm{奥秘的信德}{mystic faith}和善生的恩宠，这一切都在传述主耶稣的十字架。\par

39. 教会在她们内是美丽的。因此，天主圣言对她说：\BibleQuote{我的爱卿，你全美丽，毫无瑕疵，}因为罪过已被洗涤。\BibleQuote{你从\Place{黎巴嫩}下来，我的新娘，你从\Place{黎巴嫩}下来，从信德之始你将经过并前行，}\BibleRef{歌 4:7}因为弃绝了世界，她经历了暂世的事物而前行到基督。天主圣言又对她说：\BibleQuote{我的爱卿，你多么美丽，多么可爱，在欢愉中！你的身材犹如棕榈树，你的乳房犹如葡萄累累。}\BibleRef{歌 7:6}\par

40. 教会回答他说：\BibleQuote{谁把你交给我，像吮过我母亲乳房的兄弟？好叫我在外边遇见你时，能亲吻你，而不致被人轻视。我要引领你走进我母亲的家，进入怀孕我者的内室；你给我讲解。}\BibleRef{歌 8:1}你看，她因恩宠的恩赐而喜悦，渴望达至最深的\OriginalTerm{奥秘}{mysteries}，并将一切情感献给基督。她仍在寻求，仍在挑动他的爱，且请求\Place{耶路撒冷}的女子为她挑动他，渴望藉着她们——那忠实灵魂的美丽——使她的净配激发起对她更为丰富的爱。\par

41. 因此，主耶稣自己，被如此热切的爱和美丽优雅的可爱所邀请，因为现在受过洗的人不再有罪污，遂对教会说：\BibleQuote{请将我有如印玺，放在你的心上，有如印玺，放在你的肩上；}\BibleRef{歌 8:6}这就是说，我的爱卿，你是可爱的，你全美丽，什么也不欠缺。将我有如印玺放在你的心上，好使你的信德在圣事的圆满中闪耀。也让你的善工闪耀，显示出天主的肖像，你原是按照他的肖像受造的。不要让任何迫害减少你的爱，这爱，大水不能熄灭，江河不能淹没。\par

42. 然后记住，你们领受了圣神的印记；就是智慧与聪敏的神，超见与刚毅的神，明达与孝爱的神，以及敬畏上主的神，\BibleRef{依 11:2}并要保全你们所领受的。天主圣父封印了你们，主基督坚固了你们，并将圣神的\OriginalTerm{抵押}{earnest}放在你们心中，\BibleRef{格后 5:5}正如你们在宗徒的书信中所学到的。\par

% structure: p0063 source=h2 target=section reason=relative-heading-rank; base=h2

\section{第八章}

论主祭台的奥秘之宴。为避免有人轻视它，\Person{圣盎博罗削}指出，它比犹太人的神圣礼仪更为古老，因为在\Person{默基瑟德}的祭献中已预显了它，并且远胜\OriginalTerm{玛纳}{manna}，因为它是基督的圣体。\par

43. 被洁净的子民，富有这些装饰，急忙走向基督的祭台，说：\BibleQuote{我要走向天主的祭台，走向那使我喜乐的天主；}因为卸去了旧日错谬的皮囊，以鹰的青春更新，它急忙去接近那属天的宴席。它来到，看见圣祭台预备好了，便呼喊：\BibleQuote{你在我面前摆设了筵席。}\Person{达味}引出民众的发言说：\BibleQuote{上主是我的牧者，我什么也不缺乏。他把我放在翠绿的草地上，领我到幽静的溪水旁。}后来又说：\BibleQuote{纵使我应走过阴森的幽谷，我不怕凶险，因你与我同在。你的牧杖和短棒，是我的安慰。你在我面前，为我摆设了筵席，使我的对手观望。你用油傅了我的头，使我的杯爵满溢。}\par

44. 我们现在必须注意，免得有人看见可见之物（因为不可见之物，人的眼目不能看见也不能领悟），或许会说：\BibleQuote{天主降下玛纳，又给犹太人降下鹌鹑，}\BibleRef{出 16:13}但对他所爱的教会，他所预备的却是那些经上说：\BibleQuote{眼所未见，耳所未闻，人心所未想到的，就是天主为爱他的人所预备的一切。}\BibleRef{格前 2:9}因此，为避免有人这样说，我们将竭力证明，教会的圣事既比犹太会堂的礼仪更为古老，又比\OriginalTerm{玛纳}{manna}更为卓越。\par

45. 刚才读的\BibleBook{创世}{纪}课目显示它们更为古老，因为犹太\OriginalTerm{会堂}{synagogue}起源于\Person{梅瑟}的法律。但\Person{亚巴郎}却远为在先，他战胜敌人、救回侄子后，正享受胜利时，遇见了\Person{默基瑟德}，他带来了那些祭品，\Person{亚巴郎}恭敬地接受了。并不是\Person{亚巴郎}带来这些，而是\Person{默基瑟德}，他被描述为无父，无母，没有时日之始，也没有生命之终，好像天主子，\Person{保禄}在致希伯来人书中论他说：\BibleQuote{他永久为司祭，}他在拉丁译本中被称为正义之王与和平之王。\par

46. 你认出那是谁吗？一个自身连义人都难以做到的人，能是正义之王吗？一个连和平都难以保持的人，能是和平之王吗？他就是那按\OriginalTerm{天主性}{Godhead}而言无母的，因为他由天主父所生，与父同一性体；按\OriginalTerm{降生成人}{Incarnation}而言无父的，因为他由童贞女所生；没有起始也没有终结，因为他就是万有的起始与终结，元始与终末。因此，你们所领受的圣事，不是人的礼物，而是天主的礼物，由那祝福信德之父\Person{亚巴郎}者所带来，我们钦佩他的恩宠与作为。\par

47. 我们已经证明教会的圣事更为古老，现在当承认它们也更为优越。事实上，天主使玛纳降给祖先，并以自天降下的日用食粮养育他们，这真是一件奇妙的事；正如经上所说：\BibleQuote{人竟吃了天使的食粮。}然而，所有吃了那食粮的人，都死在旷野中；但你们所领受的食粮，那自天降下的生命之粮，却赐予永生的质素；谁若吃了这食粮，将永远不死，这食粮就是基督圣体。\par

49. 请思想，是天使的食粮更优越，还是基督的肉躯更优越？这肉躯实在是生命之体。那玛纳自天降下，这圣体却超越诸天之上；那玛纳属天，这圣体则属于诸天的主；那玛纳若存留到第二天便会朽坏，这圣体却远离一切朽坏，因为凡以圣洁之心领受的人，将不能感受到朽坏。为他们有水从磐石流出，为你们却有圣血从基督流出；水只是暂时满足他们，圣血却使你们永远饱足。犹太人喝了还会再渴，你们喝了之后却能免于口渴；那是在影像中，这却是在真理中。\par

49. 如果你们所惊奇的只是影像，那么那影像本身所预表的实体，又该是何等伟大呢！请看，在祖先身上所发生的事，确实是影像：\BibleQuote{经上说：他们喝了那跟随他们的神磐石的水，那磐石就是基督。但天主在他们大多数人身上不悦乐，因为他们都死在旷野里。这些事都是为我们的预像。}见：\BibleRef{格前 10:4}如今你们已认出，哪些更为优越，因为光明胜过影像，真理胜过预像，赐予者的圣体胜过自天降下的玛纳。\par

% structure: p0072 source=h2 target=section reason=relative-heading-rank; base=h2

\section{第九章}

为使无人因看到外形而动摇信德，此处引证了许多事例，说明外在的本性已被改变，从而证明饼成为基督的真实身体。然后，这论著以关于圣事的效果、领受者的心境等话题作结。\par

50. 或许你们会说：\Quote{我看到的是别的东西，你怎么断言我领受的是基督圣体呢？}这正是我们还需证明的一点。我们要用什么证据呢？让我们证明，这并非本性所形成之物，而是祝圣的祝福所成就的；祝福的德能大于本性的德能，因为藉着祝福，本性本身得以改变。\par

51. 梅瑟手持一根棍杖，他把它扔在地上，它就变成了蛇。见：\BibleRef{出 4:3}然后，他又抓住蛇的尾巴，它就恢复了棍杖的本性。你们看到，藉着先知职分的德能，发生了两种变化，即蛇的本性和棍杖的本性都改变了。埃及的河流本流着清纯的水；忽然间，从水源的脉管中涌出血来，没有人能喝河中的水。后来，因先知的祈祷，血止住了，水的本性又恢复了。希伯来人四面被困，一边是埃及人，一边是海；梅瑟举起他的棍杖，海水便分开，像墙壁一样坚立，在波浪间现出一条脚踪。约旦河逆流而上，反常地返回源头。见：\BibleRef{苏 3:16}这难道不明显地表明，海浪和河水的本性已被改变吗？祖先们干渴，梅瑟击打磐石，水就从磐石中流出。见：\BibleRef{出 17:6}难道恩宠没有产生一个违背本性的结果，使磐石涌出水来，而这水是它本性原本所没有的吗？玛辣是个极苦的水泉，以致干渴的人民不能饮用。梅瑟把一根木头投入水中，水就失去了苦味，恩宠瞬间调和了它。见：\BibleRef{出 15:25}在先知厄里叟的时代，有一先知子弟失落了斧头，斧头沉入水中。那丢了铁器的人求问厄里叟，厄里叟投下一块木头，铁器便浮了起来。这也显然是违背本性发生的，因为铁的本性比水更重。\par

52. 所以，我们看到，恩宠比本性具有更大的能力；然而，至今我们只谈到了先知祝福的恩宠。但若人的祝福具有改变本性的能力，那么当主和救主自己的话运作时，我们又该怎样论说那神圣的祝圣呢？你们所领受的圣事，正是因基督的圣言而成为圣事。如果厄里亚的话有那样的大能，能从天上降下火来，难道基督的圣言没有能力改变这些元素的本性吗？你们读经时，论及世界的创造，说：\BibleQuote{祂一发命，万物皆成；祂一命令，万物皆造。}难道基督的圣言，既能从虚无中造成那不存在的，反倒不能将已存在的事物，改变成它们原先不是的吗？因为赋予事物新的本性，并不比改变事物更难。\par

53. 但是，何必多用论证呢？让我们引用祂所给的实例，并以降生成人的奥迹，来证明这奥迹的真理。主耶稣由玛利亚诞生时，本性进程是否如常运行？若按常例，女人通常在与男人结合后才怀孕。然而，我们如今所造就的这个身体，正是由童贞女所生的。既然主耶稣本人并非按本性，而是由童贞女诞生，你们为何还在基督的身体上寻求本性的秩序呢？这是基督的真肉躯，曾被钉死，也被埋葬，因此，这确实是祂圣体的圣事。\par

54. 主耶稣亲自宣告：\BibleQuote{这是我的身体。}见：\BibleRef{玛 26:26}在祝福那些属天话语之前，论及的是另一种本性；在祝圣之后，所指的就是圣体。祂亲自谈论祂的圣血。在祝圣之前，它另有名称；在祝圣之后，便称为圣血。而你们回答说：阿们，意即\Quote{这是真的}。但愿内心的确信与口舌的宣认相符，但愿灵魂体会那言语所表达的真实。\par

因此，基督以这些圣事滋养祂的教会，借此使灵魂的本质得以坚固，看见她恩宠的不断增进，祂恰当地对她说：\BibleQuote{你的两胸多么美丽，我的妹妹，我的新娘！你的两胸比酒更美，你香液的芬芳压倒一切香料。我的新娘，你的嘴唇滴流纯蜜，舌下有蜜有乳；你衣服的芬芳好似黎巴嫩的馨香。我的妹妹，我的新娘，是一座锁闭的花园，是一个深锁的泉源。}祂藉此表明，奥迹应当与你们密封在一起，不被邪恶生活的行为和贞洁的玷污所破坏，不向不适合的人显露，也不因多嘴的饶舌而散布于不信者中间。因此，你们对信德的守护应当是好的，使生命的纯全与沉默得以无瑕地持守。\par

为此，教会同样守护着天乡奥迹的深奥，击退猛烈的风暴，招来春天恩宠的甘饴，深知自己的花园不会令基督不悦，于是邀请新郎说：\BibleQuote{北风啊，醒来吧！南风啊，吹向我的花园，使它的香气四溢。愿我的良人进入自己的园中，吃他树上的佳果。}因为她拥有良木佳果，它们将自己的根浸入圣泉之水，并蓬勃生长，结出美好的果实，因而如今不再被先知的斧头砍伐，反而充满福音的丰硕成果。\par

最后，上主也因它们的丰饶而喜悦，回答说：\BibleQuote{我的妹妹，我的新娘，我进入了我的花园，采了我的没药和香料，吃了我的蜂巢和蜂蜜，喝了我的酒和奶。}\BibleRef{歌 5:1} 信友们，你们要理解，为何祂讲到肉与饮料。毫无疑问，祂亲自在我们内吃喝，正如你们读到，祂说祂在我们的形象中坐监。\BibleRef{玛 25:36}\par

因此，教会看到如此浩大的恩宠，亦勉励她的子女和朋友们前来领受圣事，说：\BibleQuote{我的朋友，请吃吧！我的良人，请喝且酣醉。}\BibleRef{歌 5:1} 我们所吃所喝的是什么，圣神藉先知在别处已清楚说明，说：\BibleQuote{请你们体验，请你们观看：上主是何等的和蔼慈善！投奔祂的必获真福永欢。}在那圣事中就是基督，因为那是基督的身体，因此并非物质的食粮，而是精神的食粮。为此，宗徒论及这预象时说：\BibleQuote{我们的祖先吃了精神的食粮，饮了精神的饮品，}\BibleRef{格前 10:3} 因为天主的身体是精神的身体；基督的身体就是圣神的身体，因为圣神就是基督，如我们读到的：\BibleQuote{我们面前的圣神就是基督、上主。}\BibleRef{哀 4:20} 在伯多禄的书信中也读到：\BibleQuote{基督为我们死了。}\BibleRef{伯前 2:21} 最后，那食物坚固人心，那饮料\BibleQuote{使人心中喜乐，}正如先知所记载的。\par

因此，我们既已获得一切，便要明白我们已获得重生；但我们不要问：我们如何重生？难道我们再次进入母胎而重生吗？我在此看不出自然的常规。但在恩宠的卓越之处，并没有自然的秩序。再者，受孕的情形亦非总是出于自然常规，因为我们宣认主基督由童贞女受孕，并拒绝自然的秩序。因为玛利亚并非由男人受孕，而是因圣神有了身孕，正如玛窦所说：\BibleQuote{她因圣神有了身孕。}\BibleRef{玛 1:18} 那么，既然圣神降于童贞女身上，成就了受孕并完成了生育的工作，我们便不应怀疑，当祂降于圣洗池，或降于领受圣洗的人身上时，祂也成就了重生的真实。\par

\SourceRecord{Translated by H. de Romestin, E. de Romestin and H.T.F. Duckworth.；Nicene and Post-Nicene Fathers, Second Series；Vol. 10.；Edited by Philip Schaff and Henry Wace.；Buffalo, NY: Christian Literature Publishing Co.,；1896.；<http://www.newadvent.org/fathers/3405.htm>.}
